% ---- related_work.tex ----

Prabhāsa builds on a long tradition of research into data-efficient pretraining, morphologically-informed linguistic structure, and the role of formal grammatical frameworks in neural language modelling. This section situates the work within four overlapping research areas: the BabyLM controlled pretraining paradigm, morphology-aware masking strategies and curriculum learning, the computational treatment of Sanskrit and the Pāṇinian grammar, and the application of classical Indian logic (Navya-Nyāya) to reasoning and semantic inference. This section closes by identifying the novel synthesis that distinguishes the present approach.

The conventional scaling paradigm in language model pretraining prioritises large training budgets and extensive data corpora~\cite{Kaplan2020}. Recent work has reconsidered whether smaller models trained on limited compute can achieve reliable linguistic and structural competence through careful architectural and training choices. The BabyLM Challenge, introduced by Warstadt et al.~\cite{Warstadt2023}, establishes a controlled benchmark with a fixed 10M-token training budget and evaluates models on compositional generalisation tasks (SCAN, COGS, CFQ, BLiMP) and standard NLP benchmarks (GLUE). This framing isolates the contribution of training methodology from dataset scale and permits direct measurement of data efficiency. The 2024 iteration, reported by Hu et al.~\cite{Hu2024BabyLM}, evaluates competing approaches on the official leaderboard suite, including new subword tokenisation strategies and masking regimes.

Several noteworthy submissions to the BabyLM challenge demonstrate that structural priors can improve pretraining efficiency. LTG-BERT~\cite{Samuel2023LTG} employs language-theoretic linguistic structure combined with subword boundary masking. ELC-BERT~\cite{CharpentierSamuel2023} integrates explicit morphological parsing to inform masking. GPT-BERT~\cite{CharpentierSamuel2024} adapts causal language modelling objectives to the BabyLM regime. These submissions recover small but consistent gains on compositional tasks by encoding structural priors in the masking schedule or tokenisation. The official BabyLM baselines (Strict BLiMP 74.53 and Strict-Small BLiMP 65.08)~\cite{Hu2024BabyLM} establish reference points for competitive submissions.

\subsection{The Pāṇinian Grammatical Tradition: Kāraka Roles and Paribhāṣā Metarules}

Pāṇini's Aṣṭādhyāyī (circa 500 BCE) is a generative grammar of Sanskrit that derives surface word forms from abstract stems through roughly four thousand ordered rules (sūtras) governing morphology, phonology (sandhi), and argument structure. Kiparsky and Staal~\cite{KiparskyStaal1969} characterise this system as a derivation from a semantic deep level to a phonological surface level, with the kāraka relations occupying the deep, meaning-bearing layer; their analysis anticipates the deep-versus-surface distinction of generative grammar and grounds the link between Pāṇinian rules and compositional semantics. Cardona~\cite{Cardona1974Karaka} gives the authoritative treatment of the six kārakas (kartṛ the agent, karman the patient, karaṇa the instrument, sampradāna the recipient, apādāna the source, and adhikaraṇa the locus), showing that each is defined by the participant's relationship to the action rather than by surface case or word order. Because kārakas are semantic rather than positional, they form a language-independent target for a role-aware pretraining signal, even when transferred to English.

A distinctive feature of the Pāṇinian system is its layer of paribhāṣā: interpretive metarules that derive no forms themselves but govern how, and in what order, the substantive sūtras apply. Kielhorn's edition and translation of the Paribhāṣenduśekhara of Nāgojībhaṭṭa~\cite{Kielhorn1868Paribhasa} remains the standard reference for this metalinguistic apparatus. The adaptive-masking mechanism studied in this work is motivated by direct analogy to paribhāṣā: just as a metarule modulates the application of an underlying sūtra according to context, the masking schedule modulates the otherwise-uniform masking probability according to morpheme and kāraka-role boundaries, rather than treating masking as a context-free operation. Computationally, the Pāṇinian derivation system is realised by the Vidyut prakriyā generator~\cite{Prasad2024Vidyut}, while constraint-based and lexicon-driven parsers recover kāraka structure from running text~\cite{Kulkarni2010Parser, Goyal2016Sanskrit}. The present work uses Vidyut to derive the morpheme-boundary and role annotations that drive its N-hot embeddings and adaptive masking, and extracts kāraka annotations from the Digital Corpus of Sanskrit.

\subsection{Śābdabodha and Gadādhara's Vyutpattivāda: Navya-Nyāya Semantics}

The Navya-Nyāya (New Logic) school, formalised from the thirteenth century onward, developed a precise technical vocabulary for reference, qualification, and inference. Its theory of language centres on śābdabodha, the verbal cognition produced by understanding a sentence: meaning is reconstructed compositionally from individual word meanings together with the grammatical relations, notably the kāraka relations, that bind them. Matilal~\cite{Matilal1990WordWorld} presents this account of śābdabodha and Navya-Nyāya reference in modern logical terms. The most systematic classical treatment of how a sentence's meaning is derived is Gadādhara's Vyutpattivāda; Ganeri~\cite{Ganeri1999SemanticPowers} analyses this treatise in detail, showing how it articulates compositionality, the role of relational qualifiers, and quantifier scope within a formal semantics of testimony and knowing.

These ideas motivate two components of the present design. The supervised role-prediction auxiliary objective operationalises śābdabodha at the token level: the model is trained to recover kāraka-role structure from morpheme boundaries and verbal context, mirroring the compositional derivation of sentence meaning that the theory describes. The post-training reasoning scaffold adopts the Pañcāvayava (five-member) inference schema of the Nyāya tradition, training the model to decompose an inference into explicit, individually valid epistemic steps, grounded in the classical framework of epistemic validity rather than a contemporary fine-tuning template.

\subsection{Morphology-Aware Pretraining and Curriculum Masking}

The insight that masking schedules can encode linguistic structure motivates recent work on morphologically-informed adaptive masking. Traditional masked language modelling (MLM)~\cite{Devlin2019} masks input tokens uniformly at random, treating morphological boundaries and syntactic constituencies as irrelevant. Morpheme-aware masking strategies, by contrast, condition masking probability on linguistic structure, thereby increasing the difficulty of learning tasks that explicitly require compositional reasoning. Curriculum masking, in which masking difficulty escalates across training, has shown benefits for data-efficient pretraining. Søgaard et al.~\cite{Søgaard2021} demonstrate that curricula based on linguistic complexity can accelerate convergence. More recent work on adaptive masking schedules, as applied in several BabyLM submissions, dynamically adjusts masking probability based on token properties such as morpheme boundaries, argument structure, or dependency-parse depth. This approach trades increased computational cost during preprocessing against potential gains in sample efficiency and final model capacity on structured tasks.

Sanskrit and morphologically complex languages present natural testbeds for morpheme-aware pretraining. The Vidyut toolkit~\cite{Prasad2024Vidyut} provides morphological analysis and segmentation for Sanskrit. The Digital Corpus of Sanskrit~\cite{Goyal2012DCS} and the GRETIL collection~\cite{Germann1994GRETIL} supply large Sanskrit text corpora. Huet's Sanskrit Heritage platform~\cite{Huet2003SanskritHeritage} includes morphological generation and parsing. These resources enable direct measurement of morphological competence, including sandhi restoration, case-marking agreement, and stem-inflection pairing, on low-resource non-English linguistic systems where the utility of morphological structure is empirically unambiguous.

\subsection{Subword and Morpheme Tokenisation}

The choice of tokenisation scheme has profound implications for pretraining efficiency. Byte-pair encoding (BPE)~\cite{Sennrich2016BPE} and WordPiece are standard, but neither encodes morphological boundaries explicitly. Morpheme-aware tokenisation, in which vocabularies are constrained to true morphemes or morpheme-like units, reduces vocabulary size and enforces that learned representations reflect linguistic structure.

Botha and Blunsom~\cite{Botha2014Morphology} train segmentation models to recover morphemic boundaries and demonstrate downstream benefits for morphologically complex languages. More recently, Yaltik et al.~\cite{Yaltik2020} show that morpheme-based segmentation in multilingual NMT systems improves low-resource language performance. For Sanskrit, morpheme tokenisation is especially natural: the Pāṇinian grammar (circa 500 BCE) defines the language via morphological rules, and Vidyut implements this classical framework computationally~\cite{Prasad2024Vidyut}. Tokenising Sanskrit into morphemes directly instantiates this grammar.

\subsection{Rotary Embeddings and Modern Positional Encodings}

Transformer models require explicit position signals. Vaswani et al.'s sinusoidal embeddings~\cite{Vaswani2017} were standard in early work. More recent innovations include relative positional biases~\cite{Shaw2018Relative}, ALiBi~\cite{Press2022ALiBi}, and Rotary Position Embeddings (RoPE)~\cite{Su2021RoPE}. RoPE applies complex-number rotations in the embedding space, encoding position information in a way that respects the rotational structure of attention. RoPE has become standard in recent large models and provides strong compositional generalisation on length extrapolation tasks. The present work employs RoPE as a baseline positional encoding and evaluates whether morphological structure provides orthogonal benefits.

\subsection{Optimiser Choice and the Muon Optimizer}

The optimiser used during pretraining influences both convergence speed and final model capacity. Standard choices include Adam~\cite{Kingma2015Adam} and AdamW~\cite{Loshchilov2019AdamW}. Recent work has revisited the role of the optimiser in small-model pretraining. Jordan~\cite{Jordan2024Muon} proposes the Muon optimizer, a second-order method that scales well to modest model sizes and has shown competitive results on BabyLM leaderboards. Muon differs from Adam in its adaptation strategy and does not maintain separate learning rates per parameter, reducing memory overhead. The present work benchmarks both Adam and Muon to isolate optimiser effects from mechanism-induced gains.

\subsection{Causal versus Masked Language Modelling Objectives}

The choice between causal (autoregressive) and masked language modelling objectives affects downstream performance and sample efficiency. Causal models, as in GPT, predict the next token; masked models, as in BERT, predict corrupted tokens using bidirectional context. Devlin et al.~\cite{Devlin2019} demonstrate that masked models excel at cloze and extraction tasks, while causal models favour generation. For compositional generalisation, the picture is mixed. Clemons et al.~\cite{Clemons2023} show that causal objectives improve performance on COGS, while Warstadt et al.~\cite{Warstadt2023} find that masked models dominate the BabyLM benchmark suite. Charpentier and Samuel~\cite{CharpentierSamuel2024} directly compare GPT-style and BERT-style training on the same budget and corpora, finding that the objective choice alone is not decisive, but that curriculum and masking strategy interact strongly with objective type. The present work adheres to masked-language-modelling methodology, following BabyLM convention, but measures whether morphological masking and role-aware structures can improve performance on structured tasks.

\subsection{Compositional Generalisation and Benchmarks}

Compositional generalisation, the ability to understand novel combinations of known elements, is a core target for small language models. Benchmarks include SCAN~\cite{Lake2018SCAN}, which tests generalisation on synthetic compositional commands; COGS~\cite{Kim2020COGS}, which measures generalisation to novel argument structures in a controlled semantic domain; and CFQ~\cite{Keysers2020CFQ}, which evaluates compositional generalisation on semantic parsing to formal queries. BLiMP~\cite{Warstadt2020BLiMP} tests fine-grained grammatical competence across 67 diagnostic subsets. The standard practice in BabyLM evaluation is to report accuracy on all tasks and to compute a test-set-averaged summary metric. The present work evaluates all four benchmarks and reports both point estimates and 95 percent confidence intervals derived from bootstrap resampling across random seeds.

\subsection{Sample Efficiency and the Role of Structure}

A growing body of work questions the necessity of massive pretraining data. Dao et al.~\cite{Dao2023FlashAttention} improve inference efficiency. Hoffmann et al.~\cite{Hoffmann2022Chinchilla} study optimal model-to-data ratios. Ivgi et al.~\cite{Ivgi2023Efficient} show that pretraining on high-quality, repetitive corpora can match models trained on vast, diverse datasets. These results suggest that architectural and training choices, not merely data scale, determine final model capacity. The present work investigates whether encoding linguistic structure, via morphological tokenisation and kāraka-aware adaptive masking, can support data-efficient pretraining on modest budgets while maintaining compositional competence.

\subsection{Gap Analysis and Novelty}

Prior work on data-efficient pretraining and structural priors has generally been confined to well-resourced languages such as English, German, and Dutch. Morpheme-aware masking has been explored, but without the explicit coupling to a semantic role system (kārakas) or a classical grammatical framework (Pāṇinian rules). Classical Sanskrit provides a unique opportunity: the language's explicit morphological structure, the computational availability of the Pāṇinian grammar via Vidyut, and the semantic richness of the kāraka system combine to permit a fully principled instantiation of morphologically-aware pretraining. The integration of post-training Navya-Nyāya reasoning scaffolds, grounded in classical Indian logic rather than modern fine-tuning templates, represents a novel methodological contribution. The present work therefore synthesises several components: Pāṇinian morphological tokenisation combined with kāraka-aware adaptive masking during pretraining, evaluated on a small English model with cross-lingual transfer assessment, and extended with post-training epistemic reasoning constraints. Each component has precedent in the literature; their integration, applied to the BabyLM regime and measured against the official leaderboard suite, is not.
